% ===== 参考文献（Bibliography）：原文 p.163–166 | 译手 tr-ch3b =====
% 说明：文献条目一律保留英文原样（作者、标题、期刊/丛书、卷、年、页码均未翻译，未补 DOI、未联网核对），
%       条目顺序与编号照原书；p.166 为原书空白页（"This page intentionally left blank"）。
% [原文 p.163]
\begin{thebibliography}{99}

\bibitem[A]{A} J. E. Andersen, The Witten invariant of finite order mapping tori I, To appear in \emph{Journal f\"ur Reine und Angewandte Mathematik}.
\bibitem[AU1]{AU1} J. E. Andersen \& K. Ueno, Abelian conformal field theory and determinant bundles, Internat. J. Math. . 18, 919--993, 2007.
\bibitem[AU2]{AU2} J. E. Andersen \& K. Ueno, Geometric construction of modular functors from conformal field theory. J. Knot Theory Ramifications 16, 127--202, 2007.
\bibitem[AU3]{AU3} J. E. Andersen \& K. Ueno, Construction of the Reshetikhin-Turaev TQFT from conformal field theory, Preprint in preparation.
\bibitem[Ar]{Ar} M. Artin, Lectures on deformation of singularities \emph{Tata Institute Lecture Notes}, 54, 1976.
\bibitem[At]{At} M.F. Atiyah, On framings of 3-manifolds, \emph{Topology} 29 (1990) 1--7
\bibitem[BK]{BK} B. Bakalov; A. Jr. Kirillov, Lectures on tensor categories and modular functors, \emph{University Lecture Series}, 21. American Mathematical Society, Providence, RI, 2001.
\bibitem[Be]{Be} A. Beauville, Conformal blocks, fusion rules and the Verlinde formula, \emph{Proceedings of the Hirzebruch 65 Conference on Algebraic Geometry} (Ramat Gan, 1993), 75--96, \emph{Israel Math. Conf. Proc.}, 9, 1996.
\bibitem[BL]{BL} A. Beauville and Y. Laszlo, Conformal blocks and generalized theta functions, \emph{Comm. Math. Phys.}, 164, 385--419, 1994.
\bibitem[BS]{BS} A. A. Beilinson and V. V. Schechtman, Determinant bundles and Virasoro algebras, \emph{Comm. Math. Phys.} 118, 651--701, 1988.
\bibitem[BPZ]{BPZ} A. A. Belavin, A. M. Polyakov and A. B. Zamolodchikov, Infinite conformal symmetry in two-dimensional quantum field theory, \emph{Nucl. Phys. B}, 241, 333--380, 1984.
\bibitem[BF]{BF} D. Bernard and G. Felder, Fock representations and BRST cohomology in $\SL(2)$ current algebra, \emph{Comm. Math. Phys.} 127, 145--168, 1990.
\bibitem[Bl1]{Bl} C. Blanchet, Hecke algebras, modular categories and 3-manifolds quantum invariants, \emph{Topology} 39 no. 1, 193--223, 2000.
\bibitem[BHMV1]{BHMV1} C. Blanchet, N. Habegger, G. Masbaum, P. Vogel, Three-manifold invariants derived from the Kauffman bracket, \emph{Topology} 31 no. 4, 685--699, 1992.
\bibitem[BHMV2]{BHMV2} C. Blanchet, N. Habegger, G. Masbaum, P. Vogel, Topological quantum field theories derived from the Kauffman bracket, \emph{Topology} 34 no. 4, 883--927, 1995.
\bibitem[CL]{CL} E. A. Coddington and N. Levinson, \emph{Theory of ordinary differential equations} McGraw-Hill 1955.
\bibitem[DM]{DM} P. Deligne and D. Mumford, The irreducibility of the space of curves of given genus, \emph{Pub. Math. IHES}, 36, 75--109, 1969.
\bibitem[DV]{DV} R. Dijkgraaf and E. Verlinde, Modular invariance and the fusion algebra, \emph{Nucl. Phys. B(Proc. Suppl.)}, 5B, 87--97, 1988.
\bibitem[EO]{EO} T. Eguchi and H. Ooguri, Conformal and current algebras on a general Riemann surface, \emph{Nucl. Phys. B}, 282, 308--328, 1987.
\bibitem[Fa]{Fa} J. D. Fay, Theta functions on Riemann surfaces, \emph{Lecture Notes in Math.}, 352, \emph{Springer-Verlag}, 1973.
\bibitem[Fe]{Fe} G. Felder, BRST approach to minimal models, \emph{Nucl. Phys. B}, 317, 215--236, 1989.
\bibitem[GJ]{GJ} B. van Geemen and A.J. de Jong, On Hitchin's connection, \emph{J. AMS}, 11, 189--228, 1998.
\bibitem[GHP]{GHP} L. Gerritzen, F. Herrlich and M. van der Put, Stable $n$-pointed trees of projective lines, \emph{Nederl. Akad. Wetensch. Indag. Math.}, 50, 131--163, 1988.
\bibitem[G]{G} J. Grove, Constructing TQFTs from modular functors, \emph{J. Knot Theory Ramifications} 10 no. 8, 1085--1131, 2001.

% [原文 p.164]
\bibitem[Ha]{Ha} R. Hartshorne, \emph{Algebraic Geometry}, \emph{Springer-Verlag}, 1977.
\bibitem[Hi]{Hi} N. J. Hitchin, Flat connections and geometric quantization, \emph{Comm. Math. Phys.}, 131, 347--380, 1990.
\bibitem[Hu]{Hu} J.E. Humphreys, Introduction to Lie Algebras and Representation Theory, \emph{Springer-Verlag}, 1980.
\bibitem[Kac]{Kac} V. Kac, Infinite dimensional Lie algebras, third edition, \emph{Cambridge University Press}, 1990.
\bibitem[Kan]{Kan} K. Kanie, Conformal field theory and the braid group, \emph{Bulletin Fac. Edu. Mie Univ.} 40, 1--43, (1989).
\bibitem[KSU1]{KSU1} T. Katsura, Y. Shimizu and K. Ueno, New bosonization and conformal field theory over $\Z$, \emph{Comm. Math. Phys.} 121, 603--622, 1988.
\bibitem[KSU2]{KSU2} T. Katsura, Y. Shimizu and K. Ueno, Formal groups and conformal field theory over $\Z$, \emph{Advanced Studies in Pure Mathematics}, 19, 347--366, 1988.
\bibitem[KSU3]{KSU3} T. Katsura, Y. Shimizu and K. Ueno, Complex cobordism ring and conformal field theory over $\Z$, \emph{Math. Ann.} 291, 551--571, 1991.
\bibitem[KNTY]{KNTY} N. Kawamoto, Y. Namikawa, A. Tsuchiya and Y. Yamada, Geometric realization of conformal field theory on Riemann surfaces, \emph{Comm. Math. Phys.}, 116(1988), 247--308.
\bibitem[KZ]{KZ} V. G. Knizhnik and A. B. Zamolodchikov, Current algebra and Wess-Zumino model in two dimensions, \emph{Nucl. Phys. B} 247, 83--103, 1984.
\bibitem[Kod]{Kod} K. Kodaira, \emph{Complex manifolds and deformation of complex structures}, \emph{Springer-Verlag}, 1985.
\bibitem[Koh]{Koh} T. Kohno, Three-manifold invariants derived from conformal field theory and projective representations of modular groups, \emph{Intern. J. Modern Phys.}, 6, 1795--1805, 1992.
\bibitem[Kon]{Kon} M. Kontsevich, Rational conformal field theory and invariants of 3-manifolds, \emph{Preprint of Centre de Physique Theorique Marseille}, CPT-88/p2189, 1988.
\bibitem[KNR]{KNR} S. Kumar, M.S. Narasimhan \& A. Ramanan, Infinite Grassmannians and moduli spaces of $G$-bundles, \emph{Math. Ann.} 300, 395--423,1994.
\bibitem[Ku]{Ku} S. Kumar, Demazure character formula in arbitrary Kac-Moody setting, \emph{Invent. Math.}, 89, 395--423, 1987.
\bibitem[Kur]{Kur} G. Kuroki, Fock space representations of affine Lie algebras and integral representations in the Wess-Zumino-Witten models, \emph{Comm. Math. Phys.} 142, 511--542, 1991.
\bibitem[L]{L} Y. Laszlo, Hitchin's and WZW connections are the same, \emph{J. Differential Geometry}, 49(1998), 547--576.
\bibitem[LS]{LS} Y. Laszlo and C. Sorger, The line bundles on the stack of parabolic $G$-bundles over curves and their sections, \emph{preprint}, 1995.
\bibitem[Ma]{Ma} O. Mathieu, Formula de caract\`eres pour les alg\`ebres de Kac-Moody g\'en\'erales, \emph{Ast\'erisque}, 159--160, 1988.
\bibitem[MS1]{MS1} G. Moore and N. Seiberg, Polynomial equations for rational conformal field theories, \emph{Phys. Lett. B}, 212, 451--460, 1988.
\bibitem[MS2]{MS2} G. Moore and N. Seiberg, Classical and quantum conformal field theory, \emph{Comm. Math. Phys.} 123, 177--254, 1989.
\bibitem[Se]{Se} G. Segal, The definition of conformal field theory, \emph{Topology, geometry and quantum field theory}, 421--577, London Math. Soc. Lecture Note Ser., 308(2004).
\bibitem[Serre]{Serre} J.-P. Serre, \emph{Algebraic Groups and Class Fields}, Springer, 1988.
\bibitem[SGA7]{SGA7} \emph{Groupes de monodromie en g\'eom\'etrie alg\'ebrique, 1967--68(SGA 7)}, Lecture Notes in Math, No. 288, No.340, 1972, 1973. (http://modular.fas.harvard.edu/sga/sga/index.html).
\bibitem[So]{So} C. Sorger, \emph{La formule de Verlinde}, \emph{S\'eminaire Bourbaki, 47\`eme ann\'ee, 1994--95, no. 794}
\bibitem[Su]{Su} T. Suzuki, Finite-dimensionality of the space of conformal blocks, \emph{preprint} 1994.
\bibitem[T]{T} T. Terada, Quelques propri\'et\'es g\'eom\'etriques du domaine de $F_1$ et le groupe de tresses color\'ees, \emph{Publ. Res. Inst. Math. Sci.} 17, 95--111, 1981.
\bibitem[TK]{TK} A. Tsuchiya and Y. Kanie, Vertex operators in conformal field theory on $\PP^1$ and monodromy representations of braid group, \emph{Advanced Studies in Pure Mathematics} 16, 297--326, 1988.
\bibitem[TUY]{TUY} A. Tsuchiya, K. Ueno and Y. Yamada, Conformal field theory on universal family of stable curves with gauge symmetries, \emph{Advanced Studies in Pure Mathematics} 19, 459--566, 1989.

% [原文 p.165]
\bibitem[Tu]{Tu} V. Turaev, Quantum invariants of knots and 3-manifolds, \emph{W. de Gruyter}, 1994.
\bibitem[U1]{U1} K. Ueno, On conformal field theory, \emph{London Math. Soc. Lecture Note} 208, 283--345, 1995.
\bibitem[U2]{U2} K. Ueno, Introduction to conformal field theory with gauge symmetries, \emph{Geometry and physics (Aarhus, 1995)}, Lecture Notes in Pure and Appl. Math. 184, 603--745, Dekker, New York, 1997.
\bibitem[V]{V} E. Verlinde, Fusion rules and modular transformations in 2d conformal field theory, \emph{Nucl. Phys. B}, 300 [FS22], 360--376, 1988.
\bibitem[Wa]{Wa} K. Walker, On Witten's 3-manifold invariants, Preliminary version \# 2, Preprint 1991.
\bibitem[Wen]{Wen} H. Wenzl, Hecke algebra of type $A_N$ and subfactors, \emph{Invent. Math.}, 92, 349--383, 1988.
\bibitem[Wit]{Wit} E. Witten, Quantum field theory and the Jones polynomial, \emph{Comm. Math. Phys.} 121, 351--399, 1989.
\bibitem[WW]{WW} E.T. Whittaker \& G.N. Watson, A Course of Modern Analysis, Fourth Edition, \emph{Cambridge Univ. Press}, 1927.

% [原文 p.166]（原书此页为空白页："This page intentionally left blank"，无条目）

% [Tutor Note: (c) 文献条目按约定保留英文原样，未翻译标题、未补 DOI、未联网核对；仅排版本地化处理：期刊/书名用 \emph 排斜体，数学记号（$\SL(2)$、$\Z$、$\PP^1$、$A_N$、$G$-bundles、$n$-pointed）仍按原书用数学体。]
% [Tutor Note: (c) [AU1] 原书印 "Internat. J. Math. . 18"，卷号前的双句点照排未删（疑为原书笔误，见报告 ERRATA）。]
% [Tutor Note: (c) [KSU1]–[KSU3] 第一作者原书印 "T. Katsura"，与正文引用处（Katsura）不合，疑为 "T. Katsura" 之误；按「文献条目本身有误只记不改」，此处照原样保留，见报告 ERRATA。]
% [Tutor Note: (c) [KNR] 原书 "395--423,1994." 逗号后无空格；[At] 末尾无句点；[So] 末尾无句点。三处均照原样保留，未擅自补标点。]
% [Tutor Note: (c) [SGA7] 条目内的网址 (http://modular.fas.harvard.edu/sga/sga/index.html) 照原样保留；preamble 已载入 hyperref，故 URL 直接以正文排出即可，未另加 \url。]
% [Tutor Note: (c) [Wa] 中的 "Preliminary version # 2" 的 # 是 LaTeX 特殊字符，已转义为 \#，排版结果为 "Preliminary version # 2"。]
\end{thebibliography}
